\begin{abstract}
A one-factor model fitted to items from two correlated traits need not represent both traits equally. We characterize its population maximum likelihood solution and derive the resulting fit indices. For two equal blocks of $m$ items with loading $\lambda$, the optimal solution represents an equal composite when the trait correlation is at least $\varphi^* = (1 + (m-1)\lambda^2)/(m(2-\lambda^2))$. Below this threshold, two equally good solutions favor different traits; at zero correlation, each represents one trait alone. For this design, we establish global optimality over the unrestricted one-factor model and extend the analysis to unequal blocks and least-squares estimation. The solutions give exact population CFI, TLI, RMSEA, and SRMR, including the trait correlations at which conventional cutoffs are met. With loadings of .7 and 4 to 16 items in total, CFI reaches .95 at correlations of .78 to .84, although two traits generate the data. Numerical studies examine binary and ordinal items and finite samples. They also show that item order can change the solution returned by lavaan's default starting values. The results provide design-specific benchmarks for interpreting one-factor fit and for assessing which traits the fitted factor represents.
\par\medskip
\noindent\textbf{Keywords:} factor analysis, fit indices, model misspecification, unidimensionality, symmetry breaking
\end{abstract}

\section{Introduction}
\label{sec:introduction}

A one-factor model fitted to items from two correlated traits can represent both traits equally or primarily one of them. The optimal representation depends on the trait correlation, the loadings, the number of items, and the estimator. Global fit addresses a different property: how closely the fitted model reproduces the item correlations. Consequently, neither acceptable fit nor the decision to fit one factor establishes that the factor represents both traits equally. We characterize the population one-factor solution in a two-block design, derive its fit indices, and determine the trait correlations at which conventional cutoffs are met.

The distinction matters for scales that combine related domains, such as the PHQ-ADS, which sums the nine depression items of the PHQ-9 and the seven anxiety items of the GAD-7. In its validation, the estimated depression--anxiety correlations were .912, .865, and .653 in three trials, and the one-factor model met the CFI cutoff of .95 in the first two trials but exceeded the RMSEA cutoff of .06 in all three. Kroenke et al.\ (2016) retained the total score on the basis of additional evidence, including bifactor indices and scree plots. Such a result leaves two issues to resolve: whether the fit criteria can distinguish the two-trait structure from a single factor, and whether the fitted factor gives comparable weight to depression and anxiety.

The same issues arise when two constructs are merged into one factor whose fit is compared with that of a two-factor model to assess whether they are distinct (R\"onkk\"o \& Cho, 2022), and in unidimensional item response theory (IRT), where a single ability is estimated from items that may depend on more than one trait (Ackerman, 1989; Reckase, 1979). In these settings, CFI, RMSEA, and SRMR are often interpreted with conventional cutoffs (Hu \& Bentler, 1999), although the performance of a cutoff depends on the model and the form of misspecification (McNeish \& Wolf, 2023).

Our main result identifies when the best one-factor approximation treats the two traits equally. Suppose that each trait has $m$ indicators, all with standardized loading $\lambda$, and that errors are independent. Under normal-theory maximum likelihood (ML), the optimal solution has equal loadings on both item blocks when the trait correlation is at least
\[
\varphi^*=\frac{1+(m-1)\lambda^2}{m(2-\lambda^2)}.
\]
Below this correlation, two optimal solutions with identical fit have larger loadings on one block than on the other, because either trait can receive the larger weight. At zero correlation, the factor reproduces one block exactly and assigns zero loadings to the other. We refer to the equal-loading solution as the \emph{composite solution} and to the unequal-loading solutions as \emph{tilted solutions}; Figure~\ref{fig:1} shows both for four items per trait with loading .7. The change between them is a symmetry-breaking transition: the population is unchanged when the two item blocks are exchanged, but an optimal fitted model below the threshold is not. We establish the global minimum on both sides of the transition, rather than only comparing selected stationary solutions. As in Cudeck and Browne's (1992) analysis of covariance-model misspecification, this approximation problem is defined in the population, before sampling variability is introduced.

The minimum discrepancy also gives exact population CFI, TLI, RMSEA, and SRMR. Kenny and McCoach (2003, pp.~342--345) derived the symmetric fitted covariance and its population CFI and TLI for the two-trait design considered here. For this design, we determine when the symmetric solution is globally optimal and obtain the optimum when it is not, so that the resulting cutoff boundaries use the global minimum throughout the correlation range. For example, with loadings of .7 and 4 to 16 items in total, CFI reaches .95 when the traits correlate .78 to .84. These correlations lie above the transition, so the approximation that meets the CFI criterion here is a composite of two traits; the distinction between dimensionality, global fit, and trait representation remains necessary in both regimes.

When two equally long blocks have different loadings, the global ML solution favors the block with higher loadings, but a second local minimum can favor the other trait. We derive an exact condition for this second minimum and characterize the sensitivity to small loading differences near the symmetric threshold. Unequal block lengths can also change which trait receives greater weight. The estimator matters independently of these design differences: for the exchangeable design, the transition under unweighted least squares occurs at $1/p$, where $p$ is the total number of items, and the weighted least-squares result explains why a composite solution persists to still lower correlations in the ordinal designs examined here. Finally, the population minimum discrepancy does not depend on item order, but an optimizer's starting values can. Our finite-sample study, conducted with lavaan (Rosseel, 2012) in R (R Core Team, 2025), shows that this dependence can change both the returned loading pattern and its discrepancy; in unequal-block designs, the returned solution can favor the weaker block and fit worse than a solution obtained from another start.

This analysis relates to three strands of previous work. First, studies of unidimensional IRT often describe the fitted ability as an average or reference composite of the generating traits (Ackerman, 1989; Ansley \& Forsyth, 1985), and Strachan et al.\ (2022) reported agreement with a linear composite conjecture for nonnegative trait correlations but instability at negative correlations. These findings concern particular item structures and estimators. Our analysis identifies an explicit transition for a two-block population and shows how common cross-loadings change the relevant correlation, providing a mechanism for different one-factor representations without assuming that all multidimensional item sets have the same unidimensional approximation. Second, fit indices depend on the measurement design: higher loadings can produce poorer fit at a fixed misspecification, a result described as the reliability paradox (Hancock \& Mueller, 2011; Heene et al., 2011; McNeish et al., 2018), and more items can move CFI and RMSEA in different directions (Kenny \& McCoach, 2003; Shi et al., 2019). Groskurth et al.\ (2024) examined one-factor fits to two-factor ordinal data, showed that fit-index behavior depends on the estimator and measurement design, and discussed how omitted cross-loadings can change the meaning of fitted factors in a two-factor model despite apparently good fit. Groskurth et al.\ (2025) used a two-factor alternative to a one-factor model in their simulation-cum-ROC approach to selecting fit indices and cutoffs. The present analysis complements these studies by characterizing the minimizing covariance and its trait representation for a specified two-block population. Third, Fukasaku et al.\ (2025) developed algebraic methods for enumerating ML candidates and illustrated changes in the configuration of minima; the present population analysis gives a global classification for the exchangeable design and explicit conditions for competing minima when equally long blocks have different loadings.

Appendices~\ref{app:ml}--\ref{app:ceiling} give the proofs or their main steps, Appendix~\ref{supp:numerical} documents the numerical study, Appendix~\ref{app:binary} treats binary items, and Online Resources 1--3 give the complete proofs.
